%% kapa: exam=1 round=20 year=2009 period=1 form=A subject=economics subject_ko=경제원론 seq=17 number=57
%% source: Q-Net 감정평가사 기출문제 (cst00309) artlSeq=5024178 file=extracted/1008/5024178/2009년 제20회 감정평가사 1차시험 기출문제/1교시.hwp
%% conversion: kordoc 4.24.51 -> tools/gen1.py (auto); review=proofread-sonnet; flags=none
\begin{question}{57}
총수요곡선이 우하향하는 이유가 \uline{아닌} 것은?
\begin{choices}
  \item 물가수준이 하락하면 화폐의 실질가치가 상승하여 소비가 증가한다.
  \item 물가수준이 하락하면 화폐보유를 줄이기 때문에 이자율이 낮아지고 투자가 증가한다.
  \item 물가수준이 상승하면 이자율이 상승하여 할부로 구매하는 내구재 소비가 감소한다.
  \item 물가수준이 상승하면 가격경쟁력이 낮아져 순수출이 감소한다.
  \item 물가수준이 상승하면 화폐수요가 증가하여 재화와 서비스에 대한 지출이 증가한다.
\end{choices}
\end{question}
\answer{5}
